% =============================================================================================
\chapter{Discussion}\label{chap:discussion}

% ---------------------------------------------------------------------------------------------
\section{Justification des choix}

Le tableau~\ref{tab:choix} rassemble les principaux choix, l'alternative écartée et la raison du choix.

\begin{center}
\footnotesize
\begin{longtable}{>{\raggedright\arraybackslash}p{3.6cm}>{\raggedright\arraybackslash}p{3.2cm}>{\raggedright\arraybackslash}p{7.4cm}}
\caption{Choix de conception et justifications}\label{tab:choix}\\
\toprule
Choix & Alternative écartée & Raison \\
\midrule
\endfirsthead
\toprule
Choix & Alternative écartée & Raison \\
\midrule
\endhead
\bottomrule
\endfoot
Prêts SBA réels & données entièrement simulées & le risque doit être appris sur de vrais défauts ; données publiques et volumineuses \\
Défaut à 60 mois & perte à toute date & horizon d'un crédit d'investissement ; une perte tardive ne dépend plus du dossier \\
Séparation par année & tirage au hasard & une banque prédit le futur avec le passé \\
Test lu une seule fois & réglages vérifiés sur le test & éviter de choisir sans le vouloir ce qui réussit sur le test \\
Trois tranches de durée & durée exacte & la durée exacte est modifiée après un défaut (fuite) \\
Sans emplois ni garanties dans le modèle & toutes les variables & ce sont des critères MCDA : éviter le double comptage \\
Trois modèles, choix parmi les interprétables & choisir le plus performant & une banque doit expliquer un refus ; LightGBM mesure le prix de ce choix \\
Calibration isotonique & probabilités brutes & le MCDA utilise la probabilité ; les données d'entraînement incluent la crise \\
Pas de rééquilibrage des classes & sur- ou sous-échantillonnage & préserver des probabilités fidèles \\
Indicateurs simulés et documentés & indicateurs réels & aucune base ouverte ne relie comptes du projet et issue du prêt \\
AHP & poids posés directement & méthode reconnue qui contrôle la cohérence \\
Filtre éliminatoire & tout passer dans le classement & certaines conditions sont des couperets (ne pas pouvoir payer ses échéances) \\
Seuil de risque sur le taux de validation & taux d'entraînement & cohérence avec la période de calibration ; sinon le filtre est inopérant \\
ELECTRE III et TRI & somme pondérée & non-compensation : un risque élevé n'est pas racheté par une bonne VAN \\
Règle TRI pessimiste & règle optimiste & prudence bancaire \\
Profils absolus quand c'est possible & tertiles du lot & une catégorie doit traduire un standard, pas un rang \\
PROMETHEE II en complément & une seule méthode & vérifier que le classement ne dépend pas de la méthode \\
$\lambda$ entre 0,5 et 0,8 dans les scénarios & $\lambda$ jusqu'à 1 & $\lambda = 1$ exige l'unanimité, cas extrême montré à part \\
Arbitrage \og suivre le classement \fg{} & maximiser le nombre de projets & ne pas écarter le meilleur projet au profit de petits projets \\
Le LLM ne calcule pas & laisser le LLM analyser & un LLM peut inventer des chiffres ; chaque chiffre doit être traçable \\
Outils et vérificateur & confiance dans le LLM & chaque nombre est contrôlé automatiquement \\
Gemini gratuit & API payante & budget nul d'un projet étudiant \\
Application configurable & application figée & réutiliser la méthode sur de vraies données \\
\end{longtable}
\end{center}

% ---------------------------------------------------------------------------------------------
\section{Limites}

\paragraph{Indicateurs simulés.} C'est la limite principale : les catégories des 40 projets ne sont stables que dans
48\,\% des tirages de simulation. La démonstration valide la \emph{méthode} et la chaîne de décision, pas les catégories
de ces projets. Plusieurs paramètres de simulation sont des hypothèses non sourcées (croissance, horizon, valeur
résiduelle, bornes de dette et de garantie, centres ESG). Les 8 projets écartés pour un DSCR inférieur à 1 sont
exactement les 8 projets \og faibles \fg{} de la simulation : ce résultat découle de la construction.

\paragraph{Lien mécanique entre critères.} Le DSCR étant imposé, les flux sont proportionnels au service de la dette :
un prêt court, au service annuel élevé, a mécaniquement des flux plus forts, donc une VAN plus élevée et un délai plus
court. Comme les prêts courts sont aussi plus risqués pour le modèle, la probabilité de défaut est corrélée à la VAN et
au délai. Ce lien passe par la tranche de durée, connue à l'octroi, et non par l'issue réelle, mais il est un artefact
de la simulation.

\paragraph{Représentativité du modèle.} Le modèle a appris sur des prêts américains garantis par un organisme public ;
il ne représente pas directement un portefeuille bancaire marocain. Il sous-estime légèrement le risque de 2015-2016
(3,52\,\% prédit pour 4,25\,\% observé) en raison de la dérive entre périodes, et manque une majorité des défauts au
seuil F1.

\paragraph{Paramètres de la méthode.} Les jugements AHP ont été initialisés à partir des poids cibles : l'AHP en
vérifie surtout la cohérence, et devrait être rempli par des experts du crédit. Les seuils $q$, $p$, $v$ et certains
profils sont relatifs au lot : ajouter un projet peut les déplacer. Le classement ELECTRE III affiché est un ordre total
de lecture, qui résume les incomparabilités (signalées par les \og positions ambiguës \fg{}).

\paragraph{Agents fondés sur un LLM.} Le vérificateur contrôle les nombres et les noms de méthodes, pas le sens des
phrases : une explication peut rester approximative et demande une relecture humaine. L'offre gratuite de Gemini est
souvent saturée, ce qui allonge les réponses ; l'application reste utilisable sans LLM.

% ---------------------------------------------------------------------------------------------
\section{Écarts au cahier des charges}

\begin{itemize}
  \item \textbf{Agent extracteur} : non réalisé, nos dossiers étant des données structurées ; l'import de fichiers et
        le calcul automatique du risque en tiennent lieu.
  \item \textbf{TRI} : non calculé ; la VAN et le délai de récupération couvrent la rentabilité et la liquidité.
  \item \textbf{Durée du prêt} : remplacée par trois tranches à cause de la fuite détectée.
  \item \textbf{Seuil du filtre} : fondé sur le taux de défaut de validation plutôt que d'entraînement.
  \item \textbf{Variable de garantie} : retirée du modèle de risque pour éviter un double comptage.
\end{itemize}
Chacun de ces écarts est motivé par la qualité des résultats ou la cohérence de la méthode et est documenté.

% ---------------------------------------------------------------------------------------------
\section{Perspectives}

\begin{itemize}
  \item Appliquer l'outil à de vrais dossiers de crédit (bilans, plans de financement) pour remplacer la simulation.
  \item Réentraîner le modèle sur des données locales, ajouter des variables macroéconomiques et recalibrer
        régulièrement pour suivre la dérive.
  \item Faire remplir les comparaisons AHP par des analystes, et ajouter le TRI et une détection d'anomalies.
  \item Réaliser l'agent extracteur : un LLM qui lit des dossiers en PDF, remplit la matrice de performance en citant
        ses sources, et la soumet à la validation des analystes.
  \item Afficher le graphe de surclassement complet d'ELECTRE III et étendre l'arbitrage à plusieurs périodes.
\end{itemize}
